\begin{abstract}
We construct a finite-dimensional complex algebra whose finite-dimensional
modules have unbounded finite projective dimensions. This disproves the little
finitistic-dimension conjecture.
\end{abstract}

\section{Introduction}
\label{sec:introduction}

For a finite-dimensional algebra \(A\) over a field \(k\), its
\emph{little finitistic dimension} is
\[
\findim A
=\sup\{\pd_A N:
       N\text{ is a finitely generated left }A\text{-module},
       \ \pd_A N<\infty\}.
\]
The little finitistic-dimension conjecture asserts that this number
is finite for every such algebra.  Thus it asks for a bound on the
lengths of the finite projective resolutions over each fixed algebra,
while allowing other modules to have infinite projective dimension.
The finitistic-dimension questions were publicized by Bass in
1960 \cite{Bass1960}.  The \emph{big} finitistic dimension \(\Findim A\)
is defined by allowing all left modules in the same supremum; the little
dimension is at most the big one.

This paper gives a negative resolution of the little
finitistic-dimension conjecture.

\begin{theorem}\label{thm:main}
There is a finite-dimensional unital complex algebra \(A\) such that,
for every integer \(m\ge1\), there is a finite-dimensional left
\(A\)-module \(N_m\) with
\[
2m-2\le\pd_A N_m<\infty.
\]
In particular, \(\findim A=\infty\).
\end{theorem}

The algebra in \Cref{thm:main} is fixed before \(m\) varies.
Consequently its finitely generated modules admit terminating
projective resolutions of unbounded lengths.  This also makes its
big finitistic dimension infinite and refutes the corresponding
universal finiteness assertion for Artin algebras, since a
finite-dimensional complex algebra is an Artin algebra.

Independently, the companion~\cite[Corollary~1.3]{OpenAITachikawa2026}
obtains finite-dimensional opposite endomorphism algebras \(\Gamma_K\)
of infinite little and big left finitistic dimensions for every field
extension \(K/k\), where \(k=\mathbb F_2(q,H_1,H_2)\), using successive
cokernels in a projective-injective resolution. This is a distinct
characteristic-two example; the algebra in Theorem~\ref{thm:main}
remains the complex example constructed here.

\subsection{Background and consequences}

The equality of the little and big finitistic dimensions is a different
question from their finiteness. Huisgen-Zimmermann disproved that equality:
for every integer \(n\ge2\), she constructed a finite-dimensional monomial
algebra with little finitistic dimension \(n\) and big finitistic dimension
\(n+1\) \cite{HuisgenZimmermann1992}. Those examples have finite finitistic
dimensions and vary with \(n\); they do not give unbounded finite projective
dimensions over a single algebra. Her survey \cite{HuisgenZimmermann1995}
describes the development of the two questions.

Important positive cases for little finitistic finiteness include
finite-dimensional monomial algebras,
by Green, Kirkman, and Kuzmanovich \cite{GreenKirkmanKuzmanovich1991},
and Artin algebras with radical cube zero, by Green and
Huisgen-Zimmermann \cite{GreenHuisgenZimmermann1991}.
Igusa and Todorov introduced numerical invariants of syzygies that give
another proof of the radical-cube-zero case and establish little
finiteness for Artin algebras of representation dimension at most three
\cite[Corollaries~7 and~9]{IgusaTodorov2005}.
More recently, Cummings proved that universal little finitistic
finiteness is equivalent to the assertion that, for every
finite-dimensional algebra, finiteness holds on the left if and only
if it holds on the right \cite[Theorem~A]{Cummings2024}.
These results address bounds within classes or reductions between
finiteness assertions. Our construction instead preserves arbitrarily
long terminating processes while keeping all defining algebra data fixed.

There is also a derived-category consequence. Say that injective left
modules \emph{generate} if the smallest triangulated subcategory of the
unbounded derived category of all left modules that contains them and is
closed under arbitrary coproducts is the whole category. Rickard proved
that injective generation implies finite big finitistic dimension
\cite[Theorem~4.3]{Rickard2019}. Thus injective left \(A\)-modules do not
generate for the algebra of \Cref{thm:main}: its big left finitistic
dimension is infinite. This applies Rickard's theorem to
\(A^{\mathrm{op}}\), matching his right-module convention.

Combining this example with Cummings' auxiliary triangular algebra gives
an extreme left--right asymmetry for both dimensions and for injective
generation.

\begin{corollary}[Extreme left--right asymmetry]
\label{cor:left-right-asymmetry}
There is a finite-dimensional unital complex algebra \(\Lambda\) such that
\[
\findim \Lambda=\Findim \Lambda=\infty,
\qquad
\findim(\Lambda^{\mathrm{op}})=\Findim(\Lambda^{\mathrm{op}})=0.
\]
All four dimensions here use left modules. Thus the little and big
finitistic dimensions of \(\Lambda\) are infinite on the left and zero
on the right. Injective left \(\Lambda\)-modules do not generate the
unbounded derived category of all left \(\Lambda\)-modules, whereas
injective right \(\Lambda\)-modules generate the unbounded derived
category of all right \(\Lambda\)-modules.
\end{corollary}

The algebra \(\Lambda\) is auxiliary: the corollary makes no assertion
about the right finitistic dimensions or right injective generation of
the original algebra \(A\). Passing from \(\Lambda\) to
\(\Lambda^{\mathrm{op}}\) reverses the two sides of the corollary.
We give the proof in \Cref{sec:left-right-asymmetry}, after constructing
\(A\).
